\section*{Bab \ref{ch:b2:randomvar} --- Peubah Acak Diskret}

\begin{solution}{exo:b2:randomvar:1}
\emph{Binomial:} $X = \sum_{i=1}^n X_i$ dengan Bernoulli
$X_i$ yang \omterm{def:b2:proba:independence}{saling bebas}; $V(X_i) = \E(X_i^2) - \E(X_i)^2 = p - p^2$, dan \omterm{def:b2:randomvar:variance}{ragam}
peubah yang \omterm{def:b2:proba:independence}{saling bebas} itu menjumlah
(\cref{thm:b2:randomvar:variancerules}):
\[
\E(X) = np, \qquad V(X) = np(1-p) .
\]
\emph{Poisson:} $\E\bigl(X(X-1)\bigr) =
\sum_{k\geq2}k(k-1)e^{-\lambda}\frac{\lambda^k}{k!} = \lambda^2
e^{-\lambda}\sum_{j\geq0}\frac{\lambda^j}{j!} = \lambda^2$, jadi
\[
V(X) = \E(X^2) - \E(X)^2
= \lambda^2 + \lambda - \lambda^2 = \lambda .
\]
\emph{Geometrik} ($q = 1 - p$): dengan menurunkan
$\sum_{k\geq0}q^k = \frac{1}{1-q}$ dua kali di dalam cakramnya
(\cref{ch:b2:powerseries}), $\sum_{k\geq2}k(k-1)q^{k-2} =
\frac{2}{(1-q)^3}$, jadi
\[
\E\bigl(X(X-1)\bigr) = pq\sum_{k\geq2}k(k-1)q^{k-2}
= \frac{2q}{p^2},
\qquad
V(X) = \frac{2q}{p^2} + \frac1p - \frac{1}{p^2}
= \frac{q}{p^2} = \frac{1-p}{p^2} .
\]
\end{solution}

\begin{solution}{exo:b2:randomvar:2}
\emph{Jumlah:} untuk $n \in \N$, menurut kesalinglepasan dan kesalingbebasannya,
\[
\P(X + Y = n)
= \sum_{k=0}^n \P(X = k)\P(Y = n - k)
= e^{-(\lambda + \mu)}\frac{1}{n!}
\sum_{k=0}^n \binom nk \lambda^k\mu^{n-k}
= e^{-(\lambda+\mu)}\frac{(\lambda + \mu)^n}{n!}
\]
menurut teorema binomial: $X + Y \sim \mathcal{P}(\lambda + \mu)$.
\emph{\omterm{def:b2:randomvar:law}{Hukum} bersyarat:} untuk $0 \leq k \leq n$,
\[
\P(X = k \mid X + Y = n)
= \frac{\P(X = k)\P(Y = n - k)}{\P(X + Y = n)}
= \binom nk
\Bigl(\frac{\lambda}{\lambda+\mu}\Bigr)^{k}
\Bigl(\frac{\mu}{\lambda+\mu}\Bigr)^{n-k} ,
\]
yakni \omterm{def:b2:randomvar:law}{hukum} binomial $\mathcal{B}\bigl(n,
\frac{\lambda}{\lambda+\mu}\bigr)$: bila cacah totalnya diketahui,
setiap \omterm{def:b2:proba:space}{kejadian} secara \omterm{def:b2:proba:independence}{saling bebas} ``memilih'' sumber pertamanya
dengan peluang yang sebanding dengan lajunya.
\end{solution}

\begin{solution}{exo:b2:randomvar:3}
Ketika $k$ mainan masih kurang, setiap kotak baru membawa mainan
baru dengan peluang $\frac kn$, bebas dari masa lalunya: waktu
tunggu $W_k$ sampai mainan baru berikutnya berhukum geometrik
$\mathcal{G}\bigl(\frac kn\bigr)$, dengan $\E(W_k) = \frac nk$, dan
$T_n = W_n + W_{n-1} + \dots + W_1$ (kotak pertama selalu memberi
mainan baru: $W_n = 1$, sesuai dengan $\E = n/n$). Menurut kelinearannya,
\[
\E(T_n) = \sum_{k=1}^n \frac nk = n\sum_{k=1}^n\frac1k
\sim n\ln n ,
\]
dengan memakai $\sum_{k\leq n}\frac1k = \ln n + \gamma + o(1)$
(\cref{ch:b2:comparison}). Mengumpulkan segelintir mainan
terakhirlah yang mahal: separuh kotaknya habis untuk sisa yang sedikit itu.
\end{solution}

\begin{solution}{exo:b2:randomvar:4}
Secara \omterm{def:b2:funcseq:def}{titik demi titik}, $X(\omega) = \#\{n \geq 1 : X(\omega) \geq n\} =
\sum_{n\geq1}\mathbf{1}_{X \geq n}(\omega)$. Keluarga rangkap
$\bigl(\mathbf{1}_{X \geq n}(\omega)\,\P(\{\omega\})\bigr)_{n,
\omega}$ taknegatif, jadi Fubini untuk keluarga
(\cref{ch:b2:series}) berlaku tanpa syarat: menjumlah dahulu atas
$n$ memberi $\E(X)$, menjumlah dahulu atas $\omega$ memberi $\sum_n \P(X
\geq n)$; keduanya berhingga serentak dan sama. Untuk $X \sim
\mathcal{G}(p)$: $\P(X \geq n) = q^{n-1}$ ($q = 1-p$), jadi $\E(X) =
\sum_{n\geq1}q^{n-1} = \frac{1}{1 - q} = \frac1p$.
\end{solution}

\begin{solution}{exo:b2:randomvar:5}
\emph{Kesetangkupan:} bola ke-$i$ yang terambil merupakan bola
seragam acak dari gucinya (setiap dari $N$ bolanya sama mungkin
menempati kedudukan $i$ pada urutan pengambilannya), jadi $\P(Y_i = 1) = \frac MN =
p$ dan $\E(X) = np$ menurut kelinearannya --- tanpa perlu kesalingbebasan.

\emph{\omterm{def:b2:randomvar:variance}{Kovariansi}:} untuk $i \neq j$, $\E(Y_iY_j) = \P(\text{pengambilan }
i, j \text{ sama-sama putih}) = \frac{M(M-1)}{N(N-1)}$ (pasangan terurut
kedudukan yang berbeda memperoleh pasangan terurut bola yang berbeda,
secara seragam). Jadi
\[
\operatorname{Cov}(Y_i, Y_j)
= \frac{M(M-1)}{N(N-1)} - \frac{M^2}{N^2}
= \frac{M(N - M)}{N^2}\cdot\frac{-1}{N-1}
= -\frac{p(1-p)}{N-1} < 0 :
\]
mengambil satu bola putih membuat yang putih makin langka bagi pengambilan lainnya.
Menurut \cref{thm:b2:randomvar:variancerules},
\[
V(X) = np(1-p) + n(n-1)\Bigl(-\frac{p(1-p)}{N-1}\Bigr)
= np(1-p)\,\frac{N - n}{N - 1} \leq np(1-p) :
\]
pencuplikan tanpa pengembalian punya rerata yang sama tetapi \omterm{def:b2:randomvar:variance}{ragam} yang \emph{lebih kecil}
daripada dengan pengembalian (kesamaannya hanya untuk $n = 1$), dengan
korelasi negatifnya berperan sebagai penstabil. Untuk $n = N$
\omterm{def:b2:randomvar:variance}{ragamnya} lenyap: cacahnya lalu bersifat deterministik.
\end{solution}

\begin{solution}{exo:b2:randomvar:6}
Dengan menguraikan di sekitar $m = \E(X)$:
\[
\E\bigl((X - c)^2\bigr)
= \E\bigl((X - m)^2\bigr) + 2(m - c)\,\E(X - m) + (m - c)^2
= V(X) + (m - c)^2 ,
\]
yang minimum persis di $c = m$ dengan nilai $V(X)$ --- \omterm{def:b2:randomvar:expectation}{nilai harapan} merupakan
penduga konstan terbaik dalam rerata kuadrat.

Bila $\P(X = m) = 1$ maka $(X - m)^2$ lenyap dengan peluang
$1$, jadi $V(X) = 0$ (keluarga pendefinisinya bersuku nol kecuali pada
himpunan nol). Sebaliknya, bila $V(X) = 0$, Chebyshev
(\cref{thm:b2:randomvar:markov}) memberi $\P\bigl(\abs{X - m} \geq
\frac1n\bigr) \leq n^2\,V(X) = 0$ untuk setiap $n$; \omterm{def:b2:proba:space}{kejadian}
$\bigl\{\abs{X - m} \geq \frac1n\bigr\}$ naik menuju $\{X \neq
m\}$, jadi kekontinuan monoton (\cref{thm:b2:proba:continuity})
memberikan $\P(X \neq m) = 0$.
\end{solution}

\begin{solution}{exo:b2:randomvar:7}
$\E(S_n) = \frac n2$ dan $V(S_n) = \frac n4$.
\emph{Markov:} $\P\bigl(S_n \geq \frac{3n}4\bigr) \leq
\frac{n/2}{3n/4} = \frac23$ --- sebuah batas konstan, tak berguna untuk
$n$ yang besar.
\emph{Chebyshev:} \omterm{def:b2:proba:space}{kejadiannya} mengakibatkan $\abs{S_n - \frac n2} \geq
\frac n4$, jadi peluangnya $\leq \frac{n/4}{(n/4)^2} =
\frac4n$ --- meluruh, tetapi hanya secara suku banyak.
\emph{Chernoff:} menurut kesalingbebasannya, $\E(e^{tS_n}) =
\prod_{i=1}^n\E(e^{tX_i}) = \bigl(\frac{1 + e^t}{2}\bigr)^n$, dan
Markov yang diterapkan pada $e^{tS_n} \geq e^{3nt/4}$ memberi, untuk setiap $t >
0$,
\[
\P\Bigl(S_n \geq \frac{3n}4\Bigr)
\leq \Bigl(\frac{1 + e^t}{2}\Bigr)^n e^{-3nt/4}
= \exp\Bigl(n\bigl(\ln\tfrac{1 + e^t}{2} - \tfrac{3t}4\bigr)\Bigr).
\]
Minimumkan pangkatnya: $\frac{\dd}{\dd t}\ln\frac{1+e^t}{2} =
\frac{e^t}{1 + e^t} = \frac34$ di $e^t = 3$, yakni $t = \ln 3$,
sehingga
\[
\P\Bigl(S_n \geq \frac{3n}4\Bigr)
\leq \Bigl(\frac{4}{2}\Bigr)^n 3^{-3n/4}
= \bigl(2 \cdot 3^{-3/4}\bigr)^n \approx (0.877)^n ,
\]
yang kecil secara eksponensial. Hierarki Markov $\to$ Chebyshev $\to$
Chernoff merupakan tangga bakunya: setiap anak tangganya menerapkan Markov pada
fungsi peubahnya yang tumbuh lebih cepat.
\end{solution}

\begin{solution}{exo:b2:randomvar:8}
Menurut teorema pemindahan (\cref{thm:b2:randomvar:transfer})
yang diterapkan pada $f\bigl(\frac{S_n}{n}\bigr)$ dengan $S_n \sim
\mathcal{B}(n, x)$:
\[
\E\Bigl[f\Bigl(\frac{S_n}{n}\Bigr)\Bigr]
= \sum_{k=0}^n f\Bigl(\frac kn\Bigr)\binom nk x^k(1-x)^{n-k}
= B_nf(x) .
\]
Tetapkan $\delta > 0$ lalu pisahkan $\abs{f(S_n/n) - f(x)}$ pada \omterm{def:b2:proba:space}{kejadian} $D
= \bigl\{\abs{\frac{S_n}{n} - x} \geq \delta\bigr\}$:
di luar $D$, selisihnya paling besar sama dengan modulus kekontinuannya
$\omega_f(\delta) = \sup_{\abs{s - t}\leq\delta}\abs{f(s) -
f(t)}$; pada $D$, paling besar $2\norm f_\infty$. Dengan mengambil \omterm{def:b2:randomvar:expectation}{nilai harapannya} lalu
memakai Chebyshev dengan $V\bigl(\frac{S_n}{n}\bigr) =
\frac{x(1-x)}{n} \leq \frac{1}{4n}$:
\[
\abs{B_nf(x) - f(x)}
\leq \E\,\abs{f(S_n/n) - f(x)}
\leq \omega_f(\delta)
+ 2\norm f_\infty\,\P(D)
\leq \omega_f(\delta) + \frac{2\norm f_\infty}{4n\delta^2} .
\]
Kekontinuan seragam $f$ pada $[0, 1]$ membuat $\omega_f(\delta) \to
0$: pilih $\delta$ lalu $n$, maka $B_nf \to f$ secara \omterm{def:b2:funcseq:def}{seragam} --- itulah
teorema hampiran Weierstrass pada \cref{ch:b2:funcseq}, yang
``lema pencacahannya'' kini terkenali sebagai ketaksamaan Chebyshev
bagi \omterm{def:b2:randomvar:law}{hukum} binomialnya.
\end{solution}

\begin{solution}{exo:b2:randomvar:9}
Uraikan $S_n^4 = \sum_{i,j,k,l}X_iX_jX_kX_l$ lalu ambil
\omterm{def:b2:randomvar:expectation}{nilai harapannya}. Menurut kesalingbebasan dan pemusatannya, setiap suku yang memuat
satu indeks yang muncul tepat sekali akan lenyap ($\E(X_i) = 0$ terfaktorkan
keluar). Suku yang bertahan: $n$ suku diagonal $\E(X_i^4)$, dan
suku yang memasangkan dua pasang indeks yang sama, $\E(X_i^2X_j^2) =
\E(X_1^2)^2$ untuk $i \neq j$, yang muncul $3n(n-1)$ kali: pilih
pasangan takterurut nilainya ($\binom n2$ cara), lalu
$\frac{4!}{2!\,2!} = 6$ cara menempatkannya pada keempat lubangnya ---
$6\binom n2 = 3n(n-1)$. Jadi, dengan $\E(X_1^2)^2
\leq \E(X_1^4)$ (Jensen atau Cauchy--Schwarz),
\[
\E(S_n^4) = n\,\E(X_1^4) + 3n(n-1)\,\E(X_1^2)^2
\leq C n^2,
\qquad C = 4\,\E(X_1^4) .
\]
Markov pada orde 4:
\[
\P\Bigl(\Bigl|\frac{S_n}{n}\Bigr| \geq \varepsilon\Bigr)
= \P\bigl(S_n^4 \geq n^4\varepsilon^4\bigr)
\leq \frac{Cn^2}{n^4\varepsilon^4}
= \frac{C}{n^2\varepsilon^4} ,
\]
yakni deret yang \omterm{def:b2:series:summable}{terjumlahkan}. Menurut Borel--Cantelli~1
(\cref{thm:b2:proba:borelcantelli}), untuk setiap $j$ \omterm{def:b2:proba:space}{kejadian}
$B_j = \limsup_n\bigl\{\abs{S_n/n} \geq \frac1j\bigr\}$ berpeluang
$0$, jadi $\P\bigl(\bigcup_j B_j\bigr) = 0$ menurut kesubaditifan
\omterm{def:b2:structures:countable}{terbilang}. Pada komplemennya --- yang berpeluang $1$
--- untuk setiap $j$ ada $N$ dengan $\abs{S_n/n} < \frac1j$ bagi
semua $n \geq N$: itu persis $S_n/n \to 0$. \omterm{def:b2:randomvar:law}{Hukum} kuat bilangan
besar berlaku di bawah momen keempat; menyingkirkan hipotesis itu
(teorema Kolmogorov) merupakan pekerjaan Tahun ke-3.
\end{solution}

\begin{solution}{exo:b2:randomvar:10}
$\P(M \leq k) = \bigl(\frac k6\bigr)^2$ (kedua dadunya paling besar
$k$, secara \omterm{def:b2:proba:independence}{saling bebas}), jadi $\P(M \geq k) = 1 -
\bigl(\frac{k-1}6\bigr)^2$ dan
\[
\E(M) = \sum_{k=1}^6\P(M \geq k)
= 6 - \frac{0 + 1 + 4 + 9 + 16 + 25}{36}
= 6 - \frac{55}{36} = \frac{161}{36} \approx 4.47 ,
\]
yang cukup jauh di atas rerata $3.5$ satu dadu tunggal, sebagaimana
layaknya sebuah maksimum.
\end{solution}

\begin{solution}{exo:b2:randomvar:11}
Dengan $I_i = \mathbf 1_{\sigma(i) = i}$: $\P(\sigma(i) = i) =
\frac{(n-1)!}{n!} = \frac1n$, jadi $\E(F_n) = n\cdot\frac1n =
1$. Untuk $i \neq j$: $\P(\sigma(i) = i, \sigma(j) = j) =
\frac{(n-2)!}{n!} = \frac1{n(n-1)}$, sehingga
\[
\operatorname{Cov}(I_i, I_j) = \frac1{n(n-1)} - \frac1{n^2}
= \frac{1}{n^2(n-1)} .
\]
Menurut perkakas \omterm{def:b2:randomvar:variance}{ragamnya}
(\cref{thm:b2:randomvar:variancerules}),
\[
V(F_n) = n\cdot\frac1n\Bigl(1 - \frac1n\Bigr)
+ n(n-1)\cdot\frac1{n^2(n-1)}
= 1 - \frac1n + \frac1n = 1 .
\]
Rerata $1$, \omterm{def:b2:randomvar:variance}{ragam} $1$, tak bergantung pada $n$ --- sesuai
dengan limit Poisson pada masalah pencocokannya
(\cref{exo:b2:proba:5}).
\end{solution}

\begin{solution}{exo:b2:randomvar:12}
$T_n = \sum_{k=1}^nG_k$ dengan $G_k \sim \mathcal G(k/n)$ menyatakan
waktu sampai terlihat mainan baru ketika $k$ mainan masih kurang, dan
tahapannya \omterm{def:b2:proba:independence}{saling bebas}. Jadi
\[
V(T_n) = \sum_{k=1}^n\frac{1 - k/n}{(k/n)^2}
\leq \sum_{k=1}^n\frac{n^2}{k^2}
\leq \frac{\pi^2}6\,n^2 ,
\]
menurut \cref{ex:b2:fourier:basel}. Dengan $\E(T_n) = nH_n$, $H_n =
\sum_1^n\frac1k$ (\cref{exo:b2:randomvar:3}), Chebyshev memberi,
untuk $\varepsilon > 0$,
\[
\P\bigl(\abs{T_n - nH_n} \geq \varepsilon\,n\ln n\bigr)
\leq \frac{\pi^2n^2/6}{\varepsilon^2n^2\ln^2 n}
= \frac{\pi^2}{6\,\varepsilon^2\ln^2n}
\xrightarrow[n\to\infty]{} 0 .
\]
Karena $H_n \sim \ln n$, membaginya dengan $n\ln n$ menunjukkan
$T_n/(n\ln n) \to 1$ dalam peluang: fluktuasi
$T_n$ berorde $n$, yang dapat diabaikan terhadap reratanya $n\ln n$.
\end{solution}

\begin{solution}{pb:b2:randomvar:1}
\textbf{1.} $\E(S_n) = \frac n2$ dan Markov
(\cref{thm:b2:randomvar:markov}) memberi $\P(S_n \geq an) \leq
\frac{n/2}{an} = \frac1{2a}$. Markov hanya mengenal reratanya: ia
tak dapat membedakan peubah yang memusat di $n/2$ dari peubah yang
tersebar antara $0$ dan $n$, jadi ia menghargai ekornya seolah-olah
seluruh massanya boleh berada di sana.

\textbf{2.} Binomial setimbangnya setangkup terhadap $n/2$
($S_n$ dan $n - S_n$ berhukum sama), jadi dengan $x = n(a -
\frac12) > 0$ kedua \omterm{def:b2:proba:space}{kejadian} $\{S_n - \frac n2 \geq x\}$ dan
$\{S_n - \frac n2 \leq -x\}$ saling lepas dan sama peluangnya:
$\P(S_n \geq an) = \frac12\P(\abs{S_n - \frac n2} \geq x)$.
Chebyshev dengan $V(S_n) = \frac n4$:
\[
\P(S_n \geq an)
\leq \frac12\cdot\frac{n/4}{n^2(a - 1/2)^2}
= \frac1{8n(a - 1/2)^2},
\]
yang bernilai $\frac2n$ di $a = \frac34$.

\textbf{3.} Menurut kesalingbebasan dan teorema hasil kalinya,
$\E(\eu^{tS_n}) = \bigl(\E \eu^{tX_1}\bigr)^n = \bigl(\frac{1
+ \eu^t}2\bigr)^n$. Markov yang diterapkan pada $\eu^{tS_n}$:
\[
\P(S_n \geq an) \leq \eu^{-tan}\Bigl(\frac{1 +
\eu^t}2\Bigr)^{\!n} = \exp\Bigl(n\bigl(\ln\tfrac{1 +
\eu^t}2 - ta\bigr)\Bigr).
\]
Turunan pangkatnya terhadap $t$ adalah $\frac{\eu^t}{1 + \eu^t}
- a$, yang lenyap di $\eu^t = \frac a{1-a}$, yakni $t^* =
\ln\frac a{1-a} > 0$; di sana $\frac{1 + \eu^{t^*}}2 =
\frac1{2(1-a)}$ dan pangkatnya sama dengan
\[
n\Bigl(-\ln 2 - \ln(1-a) - a\ln\frac a{1-a}\Bigr)
= -n\bigl(\ln2 + a\ln a + (1-a)\ln(1-a)\bigr) = -n\,I(a),
\]
dengan $I(\frac12) = 0$ dan $I'(a) = \ln\frac a{1-a} > 0$ pada
$\intoo{\frac12}1$: $I(a) > 0$. Di $a = \frac34$:
$\eu^{-I(3/4)} = \frac12(\tfrac34)^{-3/4}(\tfrac14)^{-1/4} =
2\cdot3^{-3/4}$, yakni batas pada \cref{exo:b2:randomvar:7}.

\textbf{4.} Sebanyak $n + 1$ bilangan $\binom nja^j(1-a)^{n-j}$ berjumlah
$1$, dan yang terbesar adalah yang di $j = k = an$ (modus
$\mathcal B(n, a)$ di sini adalah $\floor{(n+1)a} = k$). Maksimum
dari $n + 1$ bilangan yang berjumlah $1$ paling kecil adalah
$\frac1{n+1}$:
\[
\binom nk a^k(1-a)^{n-k} \geq \frac1{n+1}
\quad\Longrightarrow\quad
\binom nk \geq \frac{a^{-an}(1-a)^{-n(1-a)}}{n+1}
= \frac{\eu^{nH(a)}}{n+1}.
\]
Jadi $\P(S_n \geq an) \geq \binom{n}{an}2^{-n} \geq
\eu^{n(H(a) - \ln2)}/(n+1) = \eu^{-nI(a)}/(n+1)$: sampai
faktor suku banyaknya $n + 1$, pangkat Chernoff itulah kebenarannya.

\textbf{5.} $n = 100$, $a = \frac34$: Markov $\frac23$;
Chebyshev $\frac2{100} = 0.02$; Chernoff
$(2\cdot3^{-3/4})^{100} = \eu^{-100\,I(3/4)} \approx
2.1\cdot10^{-6}$, terhadap nilai persisnya $2.8\cdot10^{-7}$. Pelajarannya:
setiap momen informasi membagi batasnya secara suku banyak;
adapun momen eksponensialnya mengubah \emph{sifatnya}.

\textbf{6.} $\cosh t = \sum_{k\geq0}\frac{t^{2k}}{(2k)!}$ dan
$\eu^{t^2/2} = \sum_{k\geq0}\frac{t^{2k}}{2^kk!}$; klaimnya
menyusul suku demi suku dari $(2k)! \geq 2^kk!$, yang berlaku menurut
induksi: $(2k)! = 2k(2k-1)\cdot(2k-2)! \geq 2k\cdot
2^{k-1}(k-1)! = 2^kk!\cdot(2k-1) \geq 2^kk!$.

\textbf{7.} Menurut kesalingbebasannya, $\E\bigl(\eu^{t\sum\varepsilon_i}
\bigr) = (\cosh t)^n \leq \eu^{nt^2/2}$, jadi Markov memberi
$\P(\sum\varepsilon_i \geq s) \leq \eu^{nt^2/2 - ts}$; dengan
meminimumkannya di $t = s/n$ diperoleh $\eu^{-s^2/(2n)}$.

\textbf{8.} Dengan $X_i = \frac{1 + \varepsilon_i}2$, $\widehat
p_n - \frac12 = \frac1{2n}\sum\varepsilon_i$, jadi
$\{\widehat p_n - \frac12 \geq \delta\} =
\{\sum\varepsilon_i \geq 2n\delta\}$ dan pertanyaan 7 memberi
batas $\eu^{-(2n\delta)^2/(2n)} = \eu^{-2n\delta^2}$. \omterm{def:b2:proba:space}{Kejadian}
setangkupnya berbatas sama, dari situlah faktor $2$ bagi
$\abs{\widehat p_n - \frac12} \geq \delta$.

\textbf{9.} $\E(\eu^{tX}) = \sum_x\eu^{tx}\P(X = x)$ merupakan
deret fungsi mulus dari $t$ yang turunan suku demi
sukunya terdominasi, pada setiap selang-$t$ \omterm{def:b2:metric:compact}{kompak}, oleh
$\eu^{\abs t}\P(X = x)$ (sebab $0 \leq x \leq 1$): menurut
teorema penurunan bagi deret yang konvergen normal
(\cref{thm:b2:funcseq:differentiation}) ia dua kali
\omterm{def:b2:diffcalc:differential}{terdiferensialkan}, dan aturan hasil baginya memberi $\psi' = \E_t(X)$
dan $\psi'' = \E_t(X^2) - \E_t(X)^2$, dengan $\E_t$ menyatakan
\omterm{def:b2:randomvar:expectation}{nilai harapan} bagi bobot yang ditimbang ulang
$\eu^{tx}\P(X{=}x)/\E(\eu^{tX})$ --- taknegatif, berjumlah
$1$, ditopang oleh nilai yang sama $x \in \intcc01$. \omterm{def:b2:randomvar:variance}{Ragam}
sebuah peubah bernilai-$\intcc01$ paling besar $\frac14$: menurut
\cref{exo:b2:randomvar:6}, ia sama dengan $\min_c\E_t((X - c)^2) \leq
\E_t\bigl((X - \tfrac12)^2\bigr) \leq \tfrac14$. Taylor dengan
sisa integral, memakai $\psi(0) = 0$, $\psi'(0) = p$:
\[
\psi(t) = tp + \int_0^t(t - s)\,\psi''(s)\,\dd s
\leq tp + \frac{t^2}2\cdot\frac14,
\]
yakni $\E(\eu^{t(X - p)}) \leq \eu^{t^2/8}$ untuk semua $t$ real.

\textbf{10.} Menurut kesalingbebasannya, $\E\bigl(\eu^{t(S_n -
np)}\bigr) \leq \eu^{nt^2/8}$; Markov dan pengoptimuman $t
= 4\delta$ memberi
\[
\P(\widehat p_n - p \geq \delta)
\leq \eu^{nt^2/8 - tn\delta}\Big|_{t = 4\delta}
= \eu^{-2n\delta^2};
\]
adapun penerapannya pada peubah $1 - X_i$ (yang juga di
$\intcc01$) membatasi ekor lainnya, dari situlah batas dua sisi
$2\eu^{-2n\delta^2}$.

\textbf{11.} Chebyshev hanya perlu momen kedua lalu memberi
$\frac{p(1-p)}{n\delta^2}$; Hoeffding perlu
\emph{keterbatasan} lalu memberi $2\eu^{-2n\delta^2}$. Di $p =
\frac12$, $\delta = 0.03$: batasnya adalah
$\frac{278}{n}$ (kira-kira) berbanding
$2\eu^{-0.0018n}$; keduanya berpotongan dekat $n \approx 1200$, dan setelah
itu batas eksponensialnya menang, bahkan menang telak ($n = 5000$:
$0.056$ berbanding $2.5\cdot10^{-4}$).

\textbf{12.} Menurut Hoeffding (pertanyaan 10), $\P(\abs{\widehat
p_n - p} \geq \delta) \leq 2\eu^{-2n\delta^2} \leq \alpha$ segera
setelah $2n\delta^2 \geq \ln\frac2\alpha$, yakni $n \geq
\frac{\ln(2/\alpha)}{2\delta^2}$.

\textbf{13.} $\delta = 0.03$, $\alpha = 0.05$: $n \geq
\frac{\ln 40}{2\cdot0.0009} \approx 2049.4$: $2050$ orang.
Untuk $\delta = 0.01$: $n \geq \frac{\ln40}{0.0002} \approx
18\,445$. Ukuran populasinya tak pernah muncul sebab setiap
pemilih yang tercuplik dimodelkan sebagai pengambilan Bernoulli$(p)$ yang baru:
kesulitan jajak pendapatnya terletak pada \omterm{def:b2:randomvar:variance}{ragam} sekeping koin, bukan pada besarnya
negerinya. Memotong separuh galat batasnya berharga empat kali lipat cuplikannya
--- itulah \omterm{def:b2:randomvar:law}{hukum} $1/\delta^2$.

\textbf{14.} Chebyshev: $\P(\abs{\widehat p_n - p} \geq
\delta) \leq \frac{p(1-p)}{n\delta^2} \leq
\frac1{4n\delta^2} \leq \alpha$ untuk $n \geq
\frac1{4\alpha\delta^2}$, yakni $5556$ pada tiga poin ---
sekitar $2.7$ kali kebutuhan Hoeffding. Tanpa
pengembalian, \omterm{def:b2:randomvar:variance}{ragamnya} terkalikan $\frac{N -
n}{N-1} < 1$ (\cref{exo:b2:randomvar:5}), jadi $n$ yang sama hanya
bisa lebih baik: perhitungan dengan pengembaliannyalah yang konservatif.

\textbf{15.} $\{\widehat p_n \leq \frac12\} \subseteq
\{\widehat p_n - 0.52 \leq -0.02\}$, jadi menurut batas
Hoeffding satu sisi $\P(\widehat p_n \leq \tfrac12) \leq
\eu^{-2n(0.02)^2} \leq 0.01$ segera setelah $n \geq \frac{\ln
100}{2\cdot0.0004} \approx 5756.5$: $5757$ pemilih. Biayanya
berskala seperti kuadrat kebalikan \emph{keunggulannya}, bukan kebalikan
ketelitian yang diinginkan: pemilihan yang ketat itu mahal.

\textbf{16.} Yang dipakai: cuplikannya diambil secara seragam dan
\omterm{def:b2:proba:independence}{saling bebas} dari kumpulan pemilihnya; setiap orang yang tercuplik
menjawab, dengan jujur, dan $p$ tak bergerak selama penjajakannya. Jajak
pendapat nyata melanggar ketiganya: responden yang terjangkau dan bersedia
bukanlah cuplikan seragam (bias pemilihan dan bias tanpa jawaban),
dan jawabannya bisa tak jujur atau tak stabil. Semuanya galat
\emph{bias}: ia menggeser $\E(\widehat p_n)$ menjauhi $p$
sebesar suatu jumlah yang tak bergantung pada $n$, jadi tak ada ukuran
cuplikan yang menguranginya --- matematika pada Bagian ini hanya mengendalikan
suku fluktuasinya.

\textbf{17.} Chebyshev untuk satu kelompok berukuran $m$: $\P(\abs{
\widehat p^{(i)} - p} \geq \delta) \leq \frac{1}{4m\delta^2}
\leq \frac18$ untuk $m \geq \frac2{\delta^2}$. Bila kelompok yang keliru
kurang dari $k/2$, maka lebih dari $k/2$ di antara nilai
$\widehat p^{(i)}$ terletak pada selang buka $\intoo{p -
\delta}{p + \delta}$, begitu pula mediannya; jadi
$\{\abs{M - p} \geq \delta\}$ memaksa sedikitnya $\lceil
k/2\rceil$ kekeliruan di antara $k$ kelompok yang \omterm{def:b2:proba:independence}{saling bebas}. Batas
gabungan atas $\binom k{\lceil k/2\rceil}$ himpunan kelompok
keliru yang mungkin memberi
\[
\P(\abs{M - p} \geq \delta)
\leq \binom{k}{\lceil k/2\rceil}
\Bigl(\frac18\Bigr)^{k/2}
\leq 2^k\,8^{-k/2} = 2^{-k/2} :
\]
yakni peluruhan eksponensial terhadap banyaknya kelompok, yang ditebus
dengan \omterm{def:b2:randomvar:variance}{ragam} belaka --- berguna justru ketika sukunya
takterbatas dan Hoeffding tak tersedia.

\textbf{18.} Cauchy--Schwarz
(\cref{thm:b2:randomvar:jensen}):
\[
\E(X) = \E(X\,\mathbf 1_{X>0})
\leq \sqrt{\E(X^2)}\sqrt{\E(\mathbf 1_{X>0}^2)}
= \sqrt{\E(X^2)\,\P(X > 0)} ;
\]
kuadratkan lalu bagi.

\textbf{19.} Ambil $h(a) = I(a) - 2(a - \tfrac12)^2$. Maka
$h(\tfrac12) = 0$, $h'(a) = \ln\frac a{1-a} - 4(a -
\tfrac12)$ lenyap di $\tfrac12$, dan
\[
h''(a) = \frac1a + \frac1{1-a} - 4 = \frac{1}{a(1-a)} - 4
\geq 0
\]
sebab $a(1-a) \leq \frac14$. Jadi $h'$ naik dari $0$ pada
$\intco{\frac12}1$, sehingga $h' \geq 0$ dan $h \geq 0$: $I(a)
\geq 2(a - \tfrac12)^2$.

\textbf{20.} $I(\tfrac12) = I'(\tfrac12) = 0$, $I''(a) =
\frac1{a(1-a)}$ memberi $I''(\tfrac12) = 4$, dan $I'''(\tfrac12)
= 0$ (fungsinya setangkup terhadap $\tfrac12$), jadi
$I(\tfrac12 + \delta) = 2\delta^2 + O(\delta^4)$. Pertanyaan 4
lalu membatasi ekor sejatinya \emph{dari bawah} oleh
$\eu^{-n(2\delta^2 + O(\delta^4))}/(n+1)$: untuk $\delta$ yang kecil
pangkat Hoeffding $2n\delta^2$ persis secara asimtotik
--- hanya perbaikan suku-banyak-dalam-$n$ yang mungkin.

\textbf{21.} Markov: satu momen, peluruhan $1/a$, berguna hanya sebagai
mesin di balik yang lain (pertanyaan 1 menunjukkannya datar).
Chebyshev: dua momen, peluruhan $\frac{V}{n\delta^2}$, tajam
tanpa hipotesis tambahan
(\cref{ex:b2:randomvar:chebsharp}), sekaligus perkakas terbaik pada
pertanyaan 23. Momen keempat (\cref{exo:b2:randomvar:9}):
peluruhan $C/n^2$, keterjumlahannya pas cukup untuk sebuah \omterm{def:b2:randomvar:law}{hukum} kuat.
Hoeffding: peubah terbatas, peluruhan $2\eu^{-2n\delta^2}$,
kuda beban Bagian III. Chernoff dengan laju persisnya
$I(a)$: momen eksponensial penuh, pangkat yang tak terkalahkan
(pertanyaan 4, 20), titik acuan bagi segala yang lain.

\textbf{22.} Bila koinnya setimbang: $\P(\widehat p_n > 0.525)
\leq \P(\widehat p_n - \tfrac12 \geq 0.025) \leq
\eu^{-2n(0.025)^2}$. Bila $p = 0.55$: $\P(\widehat p_n \leq
0.525) \leq \P(\widehat p_n - 0.55 \leq -0.025) \leq
\eu^{-2n(0.025)^2}$. Kedua galatnya di bawah $0.01$ ketika
$2n(0.025)^2 \geq \ln 100$, yakni $n \geq 3684.2$: $3685$
lemparan. (Membedakan hipotesis yang terpaut $2.5$ poin berharga
sama dengan menaksir sampai ketelitian $\pm2.5$ poin.)

\textbf{23.} Hoeffding: $n \geq \frac{\ln 40}{2(0.005)^2}
\approx 73\,778$. Chebyshev dengan \omterm{def:b2:randomvar:variance}{ragam} sejatinya $p(1-p)
= 0.0099$: $n \geq \frac{0.0099}{0.05\cdot(0.005)^2} =
7920$ --- sembilan kali lebih murah. Pangkat Hoeffding
$2n\delta^2$ menghargai \omterm{def:b2:randomvar:variance}{ragamnya} pada kasus terburuknya $\frac14$,
yang pesimistis secara ganjil ketika $p = 0.01$; adapun momen kedua
yang sederhana itu lebih tahu. Perkakas yang hilang adalah batas
eksponensial yang \omterm{def:b2:randomvar:variance}{sadar-ragam} (ketaksamaan Bernstein, Tahun ke-3) ---
atau, untuk \omterm{def:b2:proba:space}{kejadian} langka, hampiran Poisson yang dibuktikan pada
\cref{ch:b2:genfun}, yang bekerja pada skala nisbi yang alami.

\textbf{24.} Tetapkan $\delta > 0$: $\sum_n 2\eu^{-2n\delta^2} <
\infty$ (deret bertipe geometrik), jadi Borel--Cantelli~1
(\cref{thm:b2:proba:borelcantelli}) memberi
$\P(\abs{\widehat p_n - p} \geq \delta \text{ tak hingga
kali}) = 0$, yakni \omterm{def:b2:proba:space}{kejadian} $E_j =
\bigcup_N\bigcap_{n\geq N}\{\abs{\widehat p_n - p} <
\tfrac1j\}$ berpeluang $1$ untuk setiap $j$. Irisan \omterm{def:b2:structures:countable}{terbilang}
$\bigcap_jE_j$ tetap berpeluang $1$
(kesubaditifan pada komplemennya), dan padanya $\widehat p_n
\to p$: itulah \omterm{def:b2:randomvar:law}{hukum} kuat bilangan besar untuk lemparan koin, dengan
keterbatasan memainkan peran yang dimainkan momen keempat pada
\cref{exo:b2:randomvar:9}.

\textbf{25.} Satu momen membeli batas yang datar; dua membeli
$1/(n\delta^2)$, tak lebih (contoh ketajamannya); empat
membeli $1/n^2$, cukup untuk meneleskopkannya menjadi \omterm{def:b2:randomvar:law}{hukum} hampir pasti;
keterbatasan membeli $\eu^{-2n\delta^2}$; dan momen
eksponensial penuh membeli laju persisnya $I$, yang tak
terkalahkan oleh metode mana pun. Menjajaki $2050$ orang memadai bagi negeri mana pun sebab
fluktuasi cuplikannya dikendalikan oleh \omterm{def:b2:randomvar:variance}{ragam} koinnya,
bukan oleh besarnya populasi --- banderol $1/\delta^2$ dan
$\ln(1/\alpha)$ berlaku semesta. Teorema limit pusat pada jilid Tahun ke-3
menggantikan ketaksamaan ini, pada skala
$\sqrt n$, dengan \omterm{def:b2:randomvar:law}{hukum} limit persis berkonstanta eksplisit ---
yang mengubah setiap batas pada masalah ini menjadi
kesamaan asimtotik.

\end{solution}
